% ============================================================
% NIDA Akademik Rapor - LaTeX Stil Tanımlamaları
% ============================================================

% Liste formatı (babel'den ÖNCE yüklenmeli — babel/arabic* çakışması için)
\usepackage{enumitem}
\setlist[itemize]{noitemsep, topsep=5pt}
\setlist[enumerate]{noitemsep, topsep=5pt}

% Türkçe dil desteği
\usepackage[utf8]{inputenc}
\usepackage[T1]{fontenc}
\usepackage[turkish]{babel}

% Sayfa düzeni
\usepackage[top=2.5cm, bottom=2.5cm, left=3cm, right=2.5cm]{geometry}
\usepackage{setspace}
\onehalfspacing

% Görsel ve tablo
\usepackage{graphicx}
\usepackage{float}
\usepackage{booktabs}
\usepackage{longtable}
\usepackage{array}
\usepackage{multirow}
\usepackage{tabularx}
\usepackage{caption}
\usepackage{subcaption}

% Kod ve UML gösterimi
\usepackage{listings}
\usepackage[svgnames]{xcolor}
\usepackage{tcolorbox}
\tcbuselibrary{listings,skins,breakable}

% Matematik (amsmath must be loaded before cleveref)
\usepackage{amsmath}
\usepackage{amssymb}

% Referans ve linkler
\usepackage{hyperref}
\usepackage{url}
\usepackage[nameinlink]{cleveref}

% Başlık formatı
\usepackage{titlesec}
\titleformat{\chapter}[display]
  {\normalfont\huge\bfseries}{\chaptertitlename\ \thechapter}{20pt}{\Huge}
\titlespacing*{\chapter}{0pt}{0pt}{40pt}

% Alt başlık formatları
\titleformat{\section}
  {\normalfont\Large\bfseries}{\thesection}{1em}{}
\titleformat{\subsection}
  {\normalfont\large\bfseries}{\thesubsection}{1em}{}
\titleformat{\subsubsection}
  {\normalfont\normalsize\bfseries}{\thesubsubsection}{1em}{}

% Paragraf ayarları
\setlength{\parindent}{1.25cm}
\setlength{\parskip}{6pt}

% Sayfa numaralandırma
\usepackage{fancyhdr}
\pagestyle{fancy}
\fancyhf{}
\fancyhead[R]{\thepage}
\fancyhead[L]{\leftmark}
\renewcommand{\headrulewidth}{0.4pt}

% UML kutuları için
\newtcblisting{umlcode}{
  colback=gray!5,
  colframe=gray!75,
  listing only,
  left=5pt,
  right=5pt,
  top=5pt,
  bottom=5pt,
  fonttitle=\bfseries,
  breakable
}

% Alıntı kutusu
\newtcolorbox{quotebox}{
  colback=blue!5,
  colframe=blue!75,
  left=10pt,
  right=10pt,
  top=5pt,
  bottom=5pt,
  breakable
}

% Uyarı kutusu
\newtcolorbox{warningbox}{
  colback=orange!10,
  colframe=orange!75,
  left=10pt,
  right=10pt,
  top=5pt,
  bottom=5pt,
  title={\textbf{Önemli Uyarı}},
  breakable
}

% Kod stili
\lstset{
  basicstyle=\ttfamily\small,
  breaklines=true,
  frame=single,
  backgroundcolor=\color{gray!10},
  keywordstyle=\color{blue},
  commentstyle=\color{green!50!black},
  stringstyle=\color{red!70!black},
  numbers=left,
  numberstyle=\tiny\color{gray},
  numbersep=5pt
}

% TypeScript/JavaScript için
\lstdefinelanguage{TypeScript}{
  keywords={const, let, var, function, return, if, else, for, while, class, interface, type, export, import, from, async, await, new, this, extends, implements},
  morecomment=[l]{//},
  morecomment=[s]{/*}{*/},
  morestring=[b]',
  morestring=[b]",
  sensitive=true
}

% Hyperref ayarları
\hypersetup{
  colorlinks=true,
  linkcolor=blue!70!black,
  citecolor=green!50!black,
  urlcolor=blue!70!black,
  pdftitle={NIDA: Yapay Zekâ Destekli Çocuk Çizimi Analiz Sistemi},
  pdfauthor={Nida ÜSTÜN},
  pdfsubject={Psikolojik Danışmanlık ve Rehberlik},
  pdfkeywords={HTP, yapay zeka, çocuk çizimi, PDR}
}

% Tablo ve şekil numaralandırma
\captionsetup[table]{labelfont=bf, textfont=it}
\captionsetup[figure]{labelfont=bf, textfont=it}

% Dipnot ayarları
\usepackage[bottom]{footmisc}

% Bibliyografya stili
\usepackage{natbib}
\bibliographystyle{apalike}

% Kelime kesme ayarları (Türkçe)
\hyphenation{psi-ko-lo-jik da-nış-man-lık reh-ber-lik pro-jek-tif}
